\section{Resumen}

El presente trabajo profesional documenta el diseño, desarrollo e implementación de \textbf{SocioUnido}, una plataforma integral de \textit{software} como servicio (\textit{SaaS}) bajo un modelo de empresa a empresa (\textit{E2E}) de marca blanca, orientada a modernizar la gestión administrativa, optimizar la recaudación financiera y segurizar el control de accesos en clubes deportivos argentinos. La solución ataca de manera directa tres ineficiencias críticas del sector: la morosidad societaria silenciosa, el fraude en los ingresos por préstamo de credenciales y el desaprovechamiento de activos e instalaciones comunes ociosas.

El ecosistema de software se compone de un panel de administración web y de aplicaciones web progresivas (\textit{PWA}) adaptativas para socios y empleados de control. Desde el punto de vista tecnológico, el proyecto destaca por tres pilares de innovación: un motor analítico de aprendizaje automático (\textit{machine learning}) para la detección temprana de anomalías en la asistencia y prevención de bajas societarias; un asistente conversacional interactivo denominado ``BotIn'' implementado sobre la API de Telegram, que utiliza procesamiento de lenguaje natural (\textit{PLN}) para automatizar consultas eliminando la fricción de descarga; y un módulo de acceso inteligente (\textit{smart access}) basado en algoritmos criptográficos de contraseñas de un solo uso basadas en tiempo (\textit{TOTP}) que permite validar la identidad y el estado de deuda del socio de forma 100\% fuera de línea (\textit{offline}), tolerando el colapso sistemático de las redes de telefonía móvil en espectáculos masivos.

El desarrollo se rigió bajo la metodología ágil Scrum a lo largo de un ciclo iterativo de trece \textit{sprints}. Para garantizar los máximos estándares de calidad, ciberseguridad e interoperabilidad del código fuente, el proceso de verificación incorporó auditorías automatizadas continuas mediante un agente autónomo de aseguramiento de calidad (\textit{QA}) montado sobre la plataforma \textit{OpenClaw} y potenciado por la inteligencia artificial de Google Gemini, complementado con análisis estático de código a través de \textit{SonarCloud} y canalizaciones de integración y despliegue continuos (\textit{CI/CD}). La viabilidad comercial de la estrategia de lanzamiento al mercado, validada mediante los marcos de las 5 C y las 4 P, prioriza inicialmente a los clubes de las categorías de ascenso metropolitano y del interior del país. En conclusión, SocioUnido demuestra ser una herramienta robusta que no solo promueve el saneamiento financiero de las instituciones, sino que impulsa la inclusión tecnológica y la reducción del impacto ambiental mediante la desmaterialización de las credenciales físicas plásticas.

\section{Palabras clave}

Plataforma \textit{SaaS} (\textit{software} como servicio), clubes de fútbol, control de accesos fuera de línea, aprendizaje automático, procesamiento de lenguaje natural, algoritmo \textit{TOTP} (contraseña de un solo uso basada en tiempo), metodología Scrum, gestión deportiva, aseguramiento de calidad automatizado, marcas blancas.

\newpage

\section{Abstract}

This professional thesis documents the design, development, and implementation of \textbf{SocioUnido}, an end-to-end (\textit{E2E}) white-label Software as a Service (\textit{SaaS}) platform engineered to modernize administrative management, optimize revenue collection, and secure access control in Argentine sports clubs. The solution directly addresses three critical sector inefficiencies: undetected member delinquency, gate fraud via credential sharing, and the underutilization of shared institutional assets and facilities.

The software ecosystem integrates an administrative web dashboard and adaptive progressive web applications (\textit{PWA}) for both members and gate staff. Technologically, the project stands out for three core innovations: a predictive machine learning engine to detect attendance anomalies and mitigate member churn; a cognitive conversational assistant named ``BotIn'' deployed on the Telegram API, leveraging natural language processing (\textit{NLP}) to resolve queries and eliminate user onboarding barriers; and a cryptographic smart access module utilizing Time-Based One-Time Password (\textit{TOTP}) algorithms that mathematically validates gate entrances 100\% \textit{offline}, thereby overcoming systematic mobile network congestion during high-concurrency stadium events.

Project execution was managed under the agile Scrum framework throughout an iterative cycle of thirteen sprints. To ensure the highest standards of software quality, cybersecurity, and interoperability, the verification process incorporated automated continuous auditing using an autonomous quality assurance (\textit{QA}) agent built on the \textit{OpenClaw} sandbox environment and powered by Google Gemini artificial intelligence, complemented by static code analysis through \textit{SonarCloud} and continuous integration and deployment (\textit{CI/CD}) pipelines. The commercial viability of the Go-To-Market strategy, validated through the 5 Cs and 4 Ps frameworks, prioritizes lower-tier and regional clubs. Ultimately, SocioUnido proves to be a robust engineering solution that fosters institutional financial recovery, promotes technological inclusion, and minimizes environmental footprint through the dematerialization of physical plastic membership cards.

\section{Keywords}

SaaS platform, football clubs, offline access control, machine learning, natural language processing, TOTP algorithm, Scrum methodology, sports management, automated quality assurance, white-label software.